% Part: first-order-logic
% Chapter: sequent-calculus
% Section: proving-things-quant

\documentclass[../../../include/open-logic-section]{subfiles}

\begin{document}

\olfileid{fol}{seq}{prq}

\olsection{带量词的推导}

\begin{ex}
给出后承式 $\lexists[x][\lnot !A(x)]
\Sequent \lnot \lforall[x][!A(x)]$ 的一个 $\Log{LK}$-推导。

在处理量词时，我们必须确保不违反本征变元条件，有时这要求我们调整某些推理的进行顺序。一般来说，先处理受本征变元条件约束的规则会有所帮助（在完成的证明中，它们会出现在更靠近结论的位置）。另外，往前看，猜测一下初始后承式可能是什么样子，也是一个好办法。在我们的例子中，它必须形如 $!A(a) \Sequent !A(a)$。这意味着，当我们“反向”应用量词规则时，必须为 $\lforall$ 规则和 $\lexists$ 规则选择同一个项——我们称之为 $a$。如果我们为每个规则选择不同的项，最终会得到像 $!A(a) \Sequent !A(b)$ 这样的结果，这当然是不可推导的。

像往常一样，我们先写下：
\begin{prooftree}
\AxiomC{}
\UnaryInf$\lexists[x][\lnot !A(x)] \fCenter \lnot \lforall[x][!A(x)]$
\end{prooftree}
我们可以选择应用 \LeftR{\exists} 规则或 \RightR{\lnot} 规则。由于 \LeftR{\exists} 规则受本征变元条件约束，宜尽早处理它，所以我们先做这一步。
\begin{prooftree}
\AxiomC{}
\UnaryInf$ \lnot !A(a) \fCenter \lnot \lforall[x][!A(x)]$
\RightLabel{\LeftR{\lexists}}
\UnaryInf$ \lexists[x][\lnot !A(x)] \fCenter \lnot \lforall[x][!A(x)]$
\end{prooftree}
反向应用 \LeftR{\lnot} 和 \RightR{\lnot} 规则，我们得到
\begin{prooftree}
\AxiomC{}
\UnaryInf$\lforall[x][!A(x)] \fCenter !A(a)$
\RightLabel{\LeftR{\lnot}}
\UnaryInf$\lnot !A(a), \lforall[x][!A(x)] \fCenter $
\RightLabel{\LeftR{\Exchange}}
\UnaryInf$\lforall[x][!A(x)], \lnot !A(a) \fCenter $
\RightLabel{\RightR{\lnot}}
\UnaryInf$ \lnot !A(a) \fCenter \lnot \lforall[x] !A(x)$
\RightLabel{\LeftR{\lexists}}
\UnaryInf$ \lexists[x] \lnot !A(x) \fCenter \lnot \lforall[x] !A(x)$
\end{prooftree}
此时，我们唯一的选择是应用 \LeftR{\forall} 规则。由于该规则不受本征变元条件的限制，我们就没有这个顾虑了。记住，我们希望最终得到一个初始后承式（形如 $!A(a) \Sequent !A(a)$），所以在应用规则时，我们应该选择 $a$ 作为 $!A$ 的自变元。
\begin{prooftree}
\Axiom$!A(a) \fCenter !A(a)$
\RightLabel{\LeftR{\lforall}}
\UnaryInf$\lforall[x][!A(x)] \fCenter !A(a)$
\RightLabel{\LeftR{\lnot}}
\UnaryInf$\lnot !A(a), \lforall[x][!A(x)] \fCenter $
\RightLabel{\LeftR{\Exchange}}
\UnaryInf$\lforall[x][!A(x)], \lnot !A(a) \fCenter $
\RightLabel{\RightR{\lnot}}
\UnaryInf$ \lnot !A(a) \fCenter \lnot \lforall[x][!A(x)]$
\RightLabel{\LeftR{\lexists}}
\UnaryInf$ \lexists[x][ \lnot !A(x)] \fCenter \lnot \lforall[x][!A(x)]$
\end{prooftree}
特别是在处理量词时，此时反复检查是否违反了本征变元条件是很重要的。由于我们应用的唯一受本征变元条件约束的规则是 \LeftR{\exists}，且本征变元 $a$ 未出现在它的下后承式（即末位后承式）中，因此这是一个正确的推导。
\end{ex}

\begin{prob}
给出下列后承式的推导：
\begin{enumerate}
\item $\Sequent (\lforall[x][!A(x)] \land \lforall[y][!B(y)]) \lif
\lforall[z][(!A(z) \land !B(z))]$。
\item $\Sequent (\lexists[x][!A(x)] \lor \lexists[y][!B(y)]) \lif
\lexists[z][(!A(z) \lor !B(z))]$。
\item $\lforall[x][(!A(x) \lif !B)] \Sequent \lexists[y][!A(y)] \lif !B$。
\item $\lforall[x][\lnot !A(x)] \Sequent \lnot\lexists[x][!A(x)]$。
\item $\Sequent \lnot\lexists[x][!A(x)] \lif \lforall[x][\lnot !A(x)]$。
\item $\Sequent \lnot\lexists[x][\lforall[y][((!A(x,y) \lif \lnot
!A(y,y)) \land (\lnot !A(y,y) \lif !A(x,y)))]]$。
\end{enumerate}
\end{prob}

\begin{prob}
给出下列后承式的推导：
\begin{enumerate}
\item $\Sequent \lnot\lforall[x][!A(x)] \lif \lexists[x][\lnot!A(x)]$。
\item $(\lforall[x][!A(x)] \lif !B) \Sequent \lexists[y][(!A(y) \lif !B)]$。
\item $\Sequent \lexists[x][(!A(x) \lif \lforall[y][!A(y)])]$。
\end{enumerate}
（这些都需要用到 $\RightR{\Contraction}$ 规则。）
\end{prob}

\end{document}